\documentclass[11pt,a4paper]{article}
\usepackage[utf8]{inputenc}
\usepackage[T1]{fontenc}
\usepackage{amsmath,amssymb,amsthm}
\usepackage{booktabs}
\usepackage[margin=2.5cm]{geometry}
\usepackage{microtype}
\usepackage[colorlinks=true,linkcolor=blue!60!black,citecolor=blue!60!black,urlcolor=blue!60!black]{hyperref}
\usepackage{xcolor}
\newcommand{\kappaLinMean}{-0.468}
\newcommand{\kappaLinMin}{-0.53}
\newcommand{\kappaLinMax}{-0.40}
\newcommand{\kappaPolyMin}{-0.87}
\newcommand{\kappaPolyMax}{-0.50}
\newcommand{\kappaSDmax}{0.87}
\newcommand{\kappaSDmin}{0.25}
\newcommand{\kappaLinFirst}{-0.52}
\newcommand{\kappaLinLast}{-0.40}
\newcommand{\kappaLinTfirst}{50}
\newcommand{\kappaLinTlast}{3200}
\newcommand{\kappaFamMin}{-0.81}
\newcommand{\kappaFamMax}{-0.59}
\newcommand{\chiHLall}{189.6}
\newcommand{\kappaCLinMean}{-0.66}
\newcommand{\kappaCFamMin}{-0.66}
\newcommand{\kappaCFamMax}{-0.63}
\newcommand{\kappaCPolyMin}{-0.75}
\newcommand{\kappaCPolyMax}{-0.51}
\newcommand{\nCells}{67}
\newcommand{\nPolyCells}{60}
\newcommand{\ctrlMin}{0.77}
\newcommand{\ctrlMax}{1.24}
\newcommand{\ctrlMean}{1.009}
\newcommand{\obsPredMin}{0.904}
\newcommand{\obsPredMax}{1.028}
\newcommand{\phiMax}{0.95}
\newcommand{\phiMin}{0.28}
\newcommand{\nBadHL}{11}
\newcommand{\badHLList}{random $d=2$, $T=25$ ($5.8\sigma$); random $d=2$, $T=50$ ($4.2\sigma$); random $d=2$, $T=100$ ($3.3\sigma$); random $d=2$, $T=400$ ($2.5\sigma$); spiral $d=3$, $T=100$ ($-2.1\sigma$); spiral $d=3$, $T=400$ ($2.0\sigma$); random $d=3$, $T=25$ ($3.8\sigma$); random $d=4$, $T=25$ ($2.7\sigma$); random $d=4$, $T=800$ ($3.0\sigma$); random $d=5$, $T=25$ ($2.4\sigma$); random $d=5$, $T=50$ ($2.5\sigma$)}
\newcommand{\nUnphysical}{0}
\newcommand{\phiTwoMin}{0.28}
\newcommand{\phiTwoMax}{0.39}
\newcommand{\rhoTwo}{0.51}
\newcommand{\chiHLdge}{109.9}
\newcommand{\nDge}{44}
\newcommand{\chiHLlin}{12.3}
\newcommand{\nLin}{7}
\newcommand{\chiOneRhoDge}{158.6}
\newcommand{\tOneRays}{5,070}
\newcommand{\tOneOk}{5,070}
\newcommand{\targetN}{8,000,000}
\newcommand{\tThreeTwoNaive}{50.1}
\newcommand{\tThreeTwoPred}{49.8}
\newcommand{\tThreeTwoBH}{0.55}
\newcommand{\tThreeFixed}{4195}
\newcommand{\tEightN}{48,000,000}
\newcommand{\tEightPi}{0.925}
\newcommand{\tEightPsi}{0.706}
\newcommand{\tEightM}{0.567}
\newcommand{\tEightSlope}{0.498}
\newcommand{\sumsLinSlopeC}{-0.53}
\newcommand{\sumsLinSlopeCsd}{0.05}
\newcommand{\sumsLinSlopeCC}{-0.41}
\newcommand{\sumsLinSlopeCCsd}{0.14}
\newcommand{\sumsPolySlopeC}{-0.50}
\newcommand{\sumsPolySlopeCsd}{0.27}
\newcommand{\sumsPolySlopeCse}{0.03}
\newcommand{\sumsPolySlopeCC}{-0.49}
\newcommand{\sumsPolySlopeCCsd}{0.49}
\newcommand{\sumsPolyN}{102}
\newcommand{\sumsLinN}{8}
\newcommand{\sumsN}{110}
\newcommand{\sumsP}{60,000}
\newcommand{\sumsCorrC}{0.11}
\newcommand{\sumsCorrCC}{0.50}
\newcommand{\sumsAmean}{-0.34}
\newcommand{\sumsAsd}{0.85}
\newcommand{\sumsNamedLegend}{$t^{2}+t+41$, $4t^{2}-2t+1$, $4t^{2}+3t+17$, $4t^{2}+30t+7$, $t^{2}+1$, $t^{3}+2$}
\newcommand{\offN}{25}
\newcommand{\offDmax}{4,000,000}
\newcommand{\offMeanLow}{0.87}
\newcommand{\offMeanHigh}{0.90}
\newcommand{\offSdHigh}{1.36}
\newcommand{\offMaxAbsHigh}{4.98}
\newcommand{\offDiagGrowth}{0.91}
\newcommand{\offSumGrowth}{0.83}
\newcommand{\offOffGrowth}{0.08}
\newcommand{\rieszX}{40,000,000}
\newcommand{\rieszRes}{0.76}
\newcommand{\rieszLone}{6.0209}
\newcommand{\rieszLtwo}{10.2438}
\newcommand{\rieszLthree}{12.9881}
\newcommand{\rieszWithin}{5}
\newcommand{\rieszChance}{0.61}
\newcommand{\rieszNpred}{30}
\newcommand{\rdX}{20,000,000}
\newcommand{\rdRes}{0.83}
\newcommand{\rdZL}{0.51}
\newcommand{\rdZZ}{0.63}
\newcommand{\rdZR}{0.57}
\newcommand{\rdZRsd}{0.76}
\newcommand{\rdSeL}{0.18}
\newcommand{\rdTopWithin}{4}
\newcommand{\mfNL}{6}
\newcommand{\mfNZ}{4}
\newcommand{\mfRpred}{13.5}
\newcommand{\mfRL}{10.4}
\newcommand{\mfRZ}{2.0}
\newcommand{\mfRctrlMean}{14.1}
\newcommand{\mfRctrlMax}{17.9}
\newcommand{\mfRrandMean}{36.5}
\newcommand{\mfRrandNinetyFive}{68.8}
\newcommand{\mfPvalue}{0.997}
\newcommand{\galTwoMean}{-0.45}
\newcommand{\galTwoSd}{0.05}
\newcommand{\galTwoN}{9}
\newcommand{\galNonMean}{-0.43}
\newcommand{\galNonSd}{0.19}
\newcommand{\galNonN}{20}
\newcommand{\exN}{25}
\newcommand{\exCmin}{0.39}
\newcommand{\exCmax}{6.64}
\newcommand{\exAlpha}{0.92}
\newcommand{\exAlphaSe}{0.08}
\newcommand{\exK}{0.43}
\newcommand{\exSigmaOne}{0.9}
\newcommand{\exSigmaTwo}{12.8}
\newcommand{\exAlphaLin}{1.19}
\newcommand{\exAlphaLinSe}{0.11}
\newcommand{\exKLin}{0.33}
\newcommand{\exKLinSe}{0.06}
\newcommand{\exSigmaOneLin}{1.7}
\newcommand{\exSigmaTwoLin}{7.4}
\newcommand{\exSigmaTwoMin}{7}
\newcommand{\exCorrC}{0.13}
\newcommand{\exCorrCC}{0.72}
\newcommand{\exSlopeC}{-0.43}
\newcommand{\exSlopeCsd}{0.12}
\newcommand{\exSlopeCse}{0.02}
\newcommand{\exSlopeCC}{-0.37}
\newcommand{\exSlopeCCsd}{0.33}
\newcommand{\exDecA}{-0.34}
\newcommand{\exDecB}{-0.58}
\newcommand{\exDecC}{-0.37}
\newcommand{\exDecAsd}{0.25}
\newcommand{\exDecBsd}{0.23}
\newcommand{\exDecCsd}{0.31}
\newcommand{\exRatioSpread}{1\cdot10^{-14}}
\newcommand{\exFracHalf}{0.86}
\newcommand{\exKone}{0.43}
\newcommand{\exKtwo}{0.37}
\newcommand{\ffCases}{6}
\newcommand{\ffNmax}{7}
\newcommand{\ffOffMax}{1.83}
\newcommand{\ffOffSwing}{1.13}
\newcommand{\fftailMaxDev}{3\cdot10^{-5}}
\newcommand{\fftailBruteN}{15}
\newcommand{\fftailBruteDev}{4\cdot10^{-16}}

\usepackage{caption}
\captionsetup{font=small,labelfont=bf}
\newtheorem{theorem}{Theorem}
\newtheorem{proposition}[theorem]{Proposition}
\newtheorem{lemma}[theorem]{Lemma}
\newtheorem{conjecture}[theorem]{Conjecture}
\newtheorem{hypothesis}[theorem]{Hypothesis}
\newcommand{\Z}{\mathbb{Z}}
\newcommand{\Q}{\mathbb{Q}}
\newcommand{\F}{\mathbb{F}}
\newcommand{\A}{\mathcal A}
\newcommand{\Sig}{\mathfrak S}
\newcommand{\Off}{\mathrm{Off}}

\title{Two routes to the variance conjecture for prime values of polynomials:\\
the off-diagonal over $\Z$ and over $\F_q[u]$}
\author{Jasper Reichardt \and Claude Fable 5.1}
\date{Working document, \today}

\begin{document}
\maketitle

\begin{abstract}
Companion to \emph{The variance of the number of prime values of a polynomial}. That paper reduces the
conjecture $\sum_{h\le H}(S_f(h)-C(f)^2)=-\tfrac12C(f)\log H+O(1)$ to Hypothesis~(E): the off-diagonal
sum over pairs of distinct roots of $f$ modulo $d$ contributes no logarithm. Here we work on the two
routes to proving~(E). Over $\Z$, for quadratics, we show that~(E) is exactly an average
equidistribution statement for the roots of $x^2\equiv D\pmod d$ of the type proved by Hooley, and we
isolate the one estimate still needed (the large moduli, Section~\ref{sec:A}). Over $\F_q[u]$ we prove
that the off-diagonal vanishes identically for all moduli of degree at most $N$, so that the analogue
of the conjecture reduces to the moduli of degree $>N$ alone, where it becomes an equidistribution
statement of Katz type; we verify the resulting exact identity numerically over $\F_3$, $\F_5$ and
$\F_7$ (Section~\ref{sec:B}).
\end{abstract}

\section{Notation and the identity to be proved}

$f\in\Z[t]$ irreducible, without fixed prime divisor; $\omega_f(p)$ roots modulo $p$;
$\nu_f(d,h)=\#\{s\bmod d: d\mid f(s),\ d\mid f(s+h)\}$; $C=C(f)$; $S_f(h)$ the pair singular series;
$b(d)=\prod_{p\mid d}(p-\omega_f(p))^{-2}$, $P_p=1-\omega_f(p)^2/(p-\omega_f(p))^2$,
$W_f(d)=d\,b(d)\prod_{p\nmid d}P_p$, $a_f(d)=W_f(d)\omega_f(d)$ ($d$ squarefree). The exact identity
of the paper is
\begin{equation}\label{eq:id}
\sum_{h\le H}\bigl(S_f(h)-C^2\bigr)=-C^2\sum_{d}a_f(d)\Bigl\{\frac Hd\Bigr\}+C^2\,\Off_f(H),\qquad
\Off_f(H)=\sum_dW_f(d)\sum_{s\ne s'}\psi_d(s'-s,H),
\end{equation}
$\psi_d(m,H)=\#\{h\le H:h\equiv m\ (d)\}-H/d$, the inner sum over ordered pairs of distinct roots of $f$
modulo $d$. The diagonal is understood (Theorems~3--4 of the paper); \textbf{Hypothesis (E)} is
$\Off_f(H)=O(1)$, and its Ces\`aro form $\Off^\ast_f(H)=O(H)$ gives the conjecture with the leading term.

\section{Route A: quadratics over $\Z$ and the roots of $x^2\equiv D$}\label{sec:A}

Let $f=at^2+bt+c$, $D=b^2-4ac$, $Q(h)=a^2h^2-D$. For $p\nmid2aD$ the roots of $f$ modulo a split prime
differ by $\pm\delta_p$, $\delta_p\equiv\sqrt D/a$. For $d$ squarefree and coprime to $2aD$ with all
prime factors split, the ordered pairs of roots modulo $d$ correspond to sign patterns
$\varepsilon\in\{-1,0,1\}^{k}$ ($k$ the number of prime factors) with multiplicity $2^{k-j(\varepsilon)}$,
$j$ the number of non-zero entries, and the difference $m_\varepsilon$ satisfies
$m_\varepsilon\equiv\varepsilon_p\delta_p\pmod p$. Hence, with $d'=\prod_{\varepsilon_p\ne0}p$,
$m_\varepsilon\equiv0\pmod{d/d'}$ and $m_\varepsilon\bmod d'$ is a root of $x^2\equiv D/a^2\pmod{d'}$.

\subsection{Reduction to root discrepancies}

For $d\le H$ write $R=d\{H/d\}$; then $\psi_d(m,H)=\mathbf 1_{m\le R}-\{H/d\}$ for $1\le m\le d$, so
\begin{equation}\label{eq:disc}
\sum_{s\ne s'}\psi_d(s'-s,H)=\Delta_d\Bigl(\Bigl\{\frac Hd\Bigr\}\Bigr),\qquad
\Delta_d(\alpha):=\#\{\text{pairs}:\ m_{ss'}\le\alpha d\}-\alpha\,\omega^\circ_f(d),
\end{equation}
the discrepancy of the multiset of root differences at the point $\alpha$, with
$\omega^\circ_f(d)=\omega_f(d)^2-\omega_f(d)$. By the correspondence above, for $d$ coprime to $2aD$,
\[
\Delta_d(\alpha)=\sum_{d'\mid d,\ d'>1}2^{\,k(d)-k(d')}\Bigl(\#\{\nu\bmod d':\ \nu^2\equiv D/a^2,\ \overline{(d/d')}\,\nu\bmod d'\le\alpha d'\}-\alpha\rho(d')\Bigr),
\]
where $\rho(d')=2^{k(d')}$ is the number of roots and $\overline{(d/d')}$ the inverse of $d/d'$ modulo $d'$
(the map $\nu\mapsto\nu\,(d/d')\,\overline{(d/d')}$ identifies the classes). Thus the off-diagonal of a
quadratic is a weighted sum, over moduli $d'$, of discrepancies of the roots of a fixed quadratic
congruence, dilated by a unit and evaluated at the point $\{H/d\}$.

\begin{proposition}\label{prop:A1}
For an irreducible quadratic $f$,
\[
\Off_f(H)=\sum_{d\le H}W_f(d)\,\Delta_d\Bigl(\Bigl\{\frac Hd\Bigr\}\Bigr)+T_f(H),\qquad
T_f(H)=\sum_{d>H}W_f(d)\Bigl(\#\{\text{pairs}:\ m_{ss'}\le H\}-\frac{H\,\omega^\circ_f(d)}d\Bigr),
\]
both sums converging absolutely.
\end{proposition}

\subsection{Small moduli: Hooley's equidistribution}

The Erd\H os--Tur\'an inequality bounds the sup-discrepancy $D^\ast_{d'}$ of the roots of $x^2\equiv D/a^2$
modulo $d'$ (dilated by a unit) by $\rho(d')/K+\sum_{k\le K}k^{-1}|\sum_\nu e(k\bar u\nu/d')|$. Hooley
\cite{Hooley1963} proved, as the key lemma of his asymptotic for $\sum_{n\le x}\tau(n^2+D)$, that the Weyl
sums $\sum_{m\le x}\sum_{\nu^2\equiv D(m)}e(k\nu/m)$ have a power saving over the trivial bound
$\sum_{m\le x}\rho(m)\asymp x$, uniformly in $k$ in a suitable range, by parametrising the solutions
$(\nu,m)$ of $\nu^2-D=m\ell$ and applying Weil's bound for Kloosterman sums; the same argument gives the
statement with the unit twist $\bar u$ and with weights that are multiplicative and bounded on average.
Since $|W_f(d)|\asymp\lambda(d)/d$ with $\lambda$ multiplicative and $\lambda(p)=p/(p-4)$ on split primes,
partial summation of a bound $\sum_{d\le x}\lambda(d)D^\ast_d\ll x(\log x)^{-1-\eta}$ gives

\begin{proposition}[conditional]\label{prop:A2}
If $\sum_{d\le x}\lambda(d)\,D^\ast_d\ll x(\log x)^{-1-\eta}$ for some $\eta>0$, then
$\sum_{d\le H}W_f(d)\Delta_d(\{H/d\})=O_f(1)$.
\end{proposition}

Hooley's lemma gives this bound with a power saving for the unweighted discrepancy of the roots of
$x^2\equiv D\pmod m$ over all $m\le x$; the extension to the weights $\lambda$ (bounded by a divisor-type
function) and to the dilations by units is routine in his method, since the Kloosterman-sum bound is
uniform in the residue class. We have not rewritten his proof in this generality and record it as the
first of the two estimates needed.

\subsection{Large moduli: the hyperbola}

For $d>H$ only differences $m\le H$ count. A pair with difference $m=h$ ($1\le h\le H$) modulo $d$ exists
iff $\nu_f(d,h)\ne0$, i.e.\ iff every prime factor of $d$ divides $h\,Q(h)$; hence
\[
T_f(H)=\sum_{h\le H}\ \sum_{\substack{d>H\\ d\mid\operatorname{rad}(hQ(h))}}W_f(d)\,\nu_f(d,h)\;-\;H\sum_{d>H}\frac{W_f(d)\,\omega^\circ_f(d)}d .
\]
Both terms are of size $\asymp\log H$ (the second by the Dirichlet series
$\sum_dW_f(d)\omega_f(d)^2d^{-s}\sim c\,\zeta(s+1)^2$; the first because $\sum_{h\le H}\tau(Q(h))\asymp H\log H$),
so their difference is a genuine secondary-term problem. Writing $W_f(d)=P(1)\,d^{-1}\prod_{p\mid d}p/(p-4)$
on split moduli and $d=Q(h)/e$ with the complementary divisor $e<Q(h)/H\le a^2H+|D|$, the first term
becomes
\[
P(1)\sum_{h\le H}\frac1{Q(h)}\sum_{\substack{e\mid Q(h)\\ e<Q(h)/H}}e\;\lambda\bigl(Q(h)/e\bigr)\quad(+\text{the terms with }p\mid h),
\]
a sum over the \emph{small} divisors $e$ of the quadratic polynomial $Q(h)$: this is the structure of
Hooley's problem $\sum_{n\le x}\tau(n^2+D)$, where the divisors $d>x$ are handled through $e=(n^2+D)/d<x$
and the count $\#\{h\le H:\ e\mid Q(h)\}=H\rho(e)/e+O(\text{discrepancy})$ is again controlled by the
equidistribution of the roots of $Q$ modulo $e$. The only new feature is the weight $\lambda(Q(h)/e)$,
which depends on the cofactor; writing $\lambda=\sum_{g\mid\cdot}\kappa(g)$ with $\kappa(p)=4/(p-4)$
reduces it to sums of the same type with $e$ replaced by $eg$. We therefore expect, but have not carried
out,

\begin{proposition}[expected]\label{prop:A3}
$T_f(H)=O_f(1)$ for every irreducible quadratic $f$, by Hooley's method applied to the small divisors
of $Q(h)$ with the multiplicative weight $\lambda$.
\end{proposition}

Propositions~\ref{prop:A2} and~\ref{prop:A3} together give Hypothesis~(E) for quadratics, and hence
(with Theorem~6 of the paper) the conjecture in Ces\`aro form. Numerically, $\Off_f(H)$ is bounded for
the $25$ quadratics tested to $H=10^5$ (paper, Table~2). The work remaining is the transcription of
\cite{Hooley1963} with weights; it is not conceptual.

\section{Route B: the function field $\F_q[u]$}\label{sec:B}

Let $\A=\F_q[u]$, $q$ odd, and $f\in\A[t]$ irreducible in $t$, separable, without fixed prime divisor.
For a prime $P$ of $\A$ write $|P|=q^{\deg P}$, $\omega_f(P)$ for the number of roots of $f$ in $\A/P$,
and define $\nu_f(d,h)$, $C(f)$, $S_f(h)$, $b$, $P_P$, $W_f$, $a_f$ exactly as over $\Z$ with $|P|$ in place
of $p$; all products converge absolutely as before. ``$h\le H$'' becomes ``$h\ne0$, $\deg h<N$'', a set of
$q^N-1$ elements, and $\log H$ becomes $N$.

\begin{theorem}\label{thm:B1}
With $a_f(d)=W_f(d)\omega_f(d)$ and $P(1)=\prod_PP_P=W_f(1)$, for every $N\ge1$,
\begin{align*}
\sum_{\substack{h\ne0\\\deg h<N}}\Bigl(\frac{S_f(h)}{C^2}-1\Bigr)
&=-\sum_{1\le\deg d\le N}a_f(d)-q^N\sum_{\deg d>N}\frac{a_f(d)}{|d|}+\bigl(1-P(1)\bigr)+\mathrm{Off}_f(N),\\
\mathrm{Off}_f(N)&=\sum_{\deg d>N}W_f(d)\sum_{s\ne s'}\Bigl(\mathbf 1_{0\le\deg\tilde m<N}-q^{N-\deg d}\Bigr),
\end{align*}
$\tilde m$ the representative of $s'-s$ of degree $<\deg d$. In particular the pairs of \emph{distinct}
roots modulo $d$ contribute nothing for $\deg d\le N$: Hypothesis~(E) is empty for small moduli.
\end{theorem}

\begin{proof}
The expansion $S_f(h)/C^2=\sum_Qb(Q)c^f_Q(h)$ over squarefree monic $Q$ holds verbatim, and for $Q\ne1$,
$\sum_{h\ne0,\deg h<N}c^f_Q(h)=\sum_{d\mid Q}|d|\mu(Q/d)\omega_f(Q/d)^2\sum_{s,s'\bmod d}\psi_d(s'-s)$ with
$\psi_d(m)=\#\{h\ne0,\deg h<N:\ h\equiv m\ (d)\}-q^{N-\deg d}$ (the subtraction is allowed because
$\sum_{d\mid Q}\mu(Q/d)=0$). For $\deg d\le N$ every class modulo $d$ contains exactly $q^{N-\deg d}$
polynomials of degree $<N$, of which one is $0$ iff $d\mid m$, so $\psi_d(m)=-\mathbf 1_{d\mid m}$, and this
includes $d=1$: $\psi_1=-1$, since $h=0$ is excluded. For $\deg d>N$ a class contains at most one nonzero
$h$ of degree $<N$, so $\psi_d(m)=\mathbf 1_{0\le\deg\tilde m<N}-q^{N-\deg d}$. Regrouping by $d$ gives
$\sum_dW_f(d)\sum_{s,s'}\psi_d(s'-s)+1$, the $+1$ compensating the excluded term $Q=1$, whose value in the
regrouped sum is $\psi_1=-1$. For $\deg d\le N$ the inner sum is $-\omega_f(d)$ (the pairs with
$s\equiv s'$), giving $-P(1)-\sum_{1\le\deg d\le N}a_f(d)$; for $\deg d>N$ the diagonal pairs give
$-\omega_f(d)q^{N-\deg d}$ and the off-diagonal pairs give $\mathrm{Off}_f(N)$. (Over $\Z$ the modulus
$d=1$ drops out because $\#\{h\le H\}=H$ exactly; here it leaves the constant $1-P(1)$.)
\end{proof}

\begin{proposition}\label{prop:B2}
$P_N(d)=P(d)+O(q^{-(N-\deg d)/2}\cdot\text{polylog})$ and, with $D_f(s)=\sum_da_f(d)|d|^{-s}$,
\[
\sum_{1\le\deg d\le N}a_f(d)=\frac N{C(f)}+c_f+O(q^{-N/2}),
\]
because, with $T=q^{-s}$, $D_f(s)=\zeta_\A(s+1)L(s+1,\chi_D)E_f(s)$ where $\zeta_\A(s+1)=1/(1-T)$,
$L(s+1,\chi_D)$ is a polynomial in $T$, and $E_f(s)=\tilde H_f(s)(1-q^{-1}T^2)^2/L(2s+2,\chi_D)$ with
$\tilde H_f$ absolutely convergent for $|T|<q$ and the zeros of the polynomial $L(2s+2,\chi_D)$ on
$|T|=q^{3/4}$ by Weil; so $D_f$ is meromorphic in $|T|<q^{3/4}$ with the single simple pole $T=1$, at which
$(1-T)D_f\to1/C(f)$ by the same local computation as over $\Z$. (The error term is in fact
$O(q^{-(3/4-\varepsilon)N})$.)
\end{proposition}

Thus, over $\F_q[u]$, the function-field form of the conjecture,
\begin{equation}\label{eq:ffconj}
\sum_{\substack{h\ne0\\\deg h<N}}\bigl(S_f(h)-C(f)^2\bigr)=-C(f)\,N+O_f(1),
\end{equation}
is equivalent to $\mathrm{Off}_f(N)=O(1)$: the whole content of Hypothesis~(E) sits in the off-diagonal
pairs modulo the moduli of degree $>N$ (the diagonal tail $q^N\sum_{\deg d>N}a_f(d)/|d|$ tends to
$1/(C(f)(q-1))$). The coefficient is $-C$ rather than $-\tfrac12C$ because the fractional parts $\{H/d\}$,
whose mean $\tfrac12$ produces the $\tfrac12$ over $\Z$, are replaced by the exact value $1$: a modulus of
degree $\le N$ misses exactly the class of $0$.

\subsection{The remainder as a finite degree-distribution problem}

For $\deg d>N$ a residue class $m$ modulo $d$ contains at most one nonzero $h$ of degree $<N$, namely its
representative $\tilde m$ when $0\le\deg\tilde m<N$. For a quadratic $f=t^2-D$ this has a strong
consequence: if $(s,s')$, $s\ne s'$, has $0\le\deg\tilde m<N$, split $d=d_1d_2$ with
$d_2=\prod\{P\mid d:\ s\equiv s'\ (P)\}$; then $d_2\mid\tilde m$, so $\deg d_2\le N-1$, and for $P\mid d_1$
the roots are $\pm\sqrt D$, so $d_1\mid\tilde m^2-4D\ne0$ and $\deg d_1\le\max(2N-2,\deg D)$. Hence
\[
\mathrm{Off}_f(N)=\sum_{N<\deg d\le M(N)}W_f(d)A_N(d)-q^N\sum_{\deg d>N}\frac{W_f(d)(\omega_f(d)^2-\omega_f(d))}{|d|},
\qquad M(N)=N-1+\max(2N-2,\deg D),
\]
with $A_N(d)$ the number of ordered pairs of distinct roots whose difference has degree $<N$; the series
is computable from $\sum_dW_f(d)\omega_f(d)^2/|d|=1$ and $\sum_dW_f(d)\omega_f(d)/|d|=\prod_P(1+b_P(\omega_P-\omega_P^2))$
(both by $\sum_{d\mid Q}\mu(Q/d)=0$). So the whole question is whether, over the moduli of degree between
$N$ and about $3N$, the differences $\pm2\sqrt D$ of the roots modulo $d$ fall among the residues of degree
$<N$ with frequency $q^{N-\deg d}$ up to a bounded total discrepancy. For a prime modulus $P$ of degree
$k$, writing the indicator of $\deg\tilde x<N$ through the additive characters trivial on $\A_{<N}$,
\[
\mathbf 1_{\deg\tilde x<N}=q^{N-k}\sum_{\deg a<k-N}\psi\Bigl(\frac{ax}{P}\Bigr),
\]
the count becomes a sum over $a\ne0$ of the Weyl sums $\sum_{\deg P=k}\psi(2a\sqrt D/P)$ of the roots of a
quadratic congruence over the primes of degree $k$: the function-field analogue of the sums of Duke,
Friedlander and Iwaniec~\cite{DFI1995}. Over $\F_q[u]$ such sums are governed by the monodromy of the
relevant Kloosterman-type sheaf and Deligne's equidistribution theorem~\cite{Katz1988}: for fixed $N,k$
and $q\to\infty$ each has square-root cancellation, $O(q^{k/2})$ against the trivial $q^k/k$, and since
only the finitely many degrees $k\le M(N)$ occur, the remainder is $O(q^{-1/2})$ in that regime.

\begin{proposition}[expected, $q\to\infty$]\label{prop:B3}
For a fixed irreducible quadratic $f=t^2-D(u)$ and fixed $N$, $\mathrm{Off}_f(N)=O_{N,\deg D}(q^{-1/2})$;
hence \eqref{eq:ffconj} holds with the $O(1)$ replaced by $O(q^{-1/2})$ as $q\to\infty$.
\end{proposition}

This is the function-field theorem to write out: its only input beyond Theorem~\ref{thm:B1} is the
equidistribution of the roots of a quadratic congruence over primes of a given degree, which in the
large-$q$ regime is a standard consequence of Deligne's theorem, in the same way as the results of
Bary-Soroker and Entin~\cite{BarySoroker2012,Entin2016} on prime values of polynomials over $\F_q[u]$.

\subsection{Exact computation}

Theorem~\ref{thm:B1} and the finite-support formula can be checked exactly, and in two independent
ways: from the $S_f(h)$ themselves (each $S_f(h)/(C^2P(1))$ is a finite product over the primes dividing
$h(h^2-4D)$; \texttt{thesis/ff.ts}) and from the roots of $f$ modulo every squarefree $d$ of degree
$\le M(N)$ (\texttt{thesis/ff-tail.ts}, with a brute-force recomputation of the finite sum for the
smallest cases). Table~\ref{tab:ff} reports the result for $q=3,5,7$ and six choices of $D$ (the cases
$D$ and $D(u+c)$ coincide, since $u\mapsto u+c$ is a degree-preserving automorphism, so only one
representative of each translation class is listed). The two computations of $\mathrm{Off}_f(N)$ agree to
$\fftailMaxDev$; the increments of $\sum_{\deg d\le N}a_f(d)$ converge to $1/C(f)$ as
Proposition~\ref{prop:B2} predicts, after a transient for $q=3$ caused by the primes of norm $3$ with
$\omega_f(P)=2$, where $P_P=1-4/(3-2)^2=-3$ and $P(1)$ is negative; the diagonal tail converges to
$1/(C(f)(q-1))$; and $\mathrm{Off}_f(N)$ is bounded, $|\mathrm{Off}_f(N)|\le\ffOffMax$ with a variation
over $N$ of at most $\ffOffSwing$, with no sign of growth, which is what \eqref{eq:ffconj} requires.

\begin{table}[h]
\centering\footnotesize\setlength{\tabcolsep}{4pt}
\caption{Theorem~\ref{thm:B1} for $f=t^2-D$ over $\F_q[u]$: LHS, $\sum_{1\le\deg d\le N}a_f(d)$ and its
increment (limit $1/C(f)$), the diagonal tail $q^N\sum_{\deg d>N}a_f(d)/|d|$ (limit $1/(C(f)(q-1))$),
and $\mathrm{Off}_f(N)$.}\label{tab:ff}
\begin{tabular}{rlrrrrrrrrr}
\toprule
$q$ & $D$ & $C(f)$ & $1-P(1)$ & $N$ & $\#h$ & LHS & $\sum a_f$ & incr. & diag.\ tail & $\mathrm{Off}_f(N)$\\
\midrule
3 & $u$ & 0.851 & 2.999 & 1 & 2 & -2.000 & 1.999 & -- & 1.168 & -1.832 \\
 & & & & 2 & 8 & -3.202 & 5.198 & 3.199 & 0.305 & -0.699 \\
 & & & & 3 & 26 & -4.758 & 5.302 & 0.104 & 0.810 & -1.645 \\
 & & & & 4 & 80 & -5.855 & 7.181 & 1.879 & 0.551 & -1.122 \\
 & & & & 5 & 242 & -7.167 & 8.308 & 1.126 & 0.527 & -1.331 \\
 & & & & 6 & 728 & -8.201 & 9.252 & 0.945 & 0.637 & -1.311 \\
 & & & & 7 & 2186 & -9.548 & 10.587 & 1.334 & 0.576 & -1.385 \\
\addlinespace
3 & $u^2+1$ & 1.459 & 3.857 & 1 & 2 & -2.000 & 5.714 & -- & -0.094 & -0.237 \\
 & & & & 2 & 8 & -2.286 & 5.306 & -0.408 & 0.127 & -0.710 \\
 & & & & 3 & 26 & -2.921 & 5.129 & -0.177 & 0.558 & -1.092 \\
 & & & & 4 & 80 & -3.530 & 6.522 & 1.394 & 0.279 & -0.586 \\
 & & & & 5 & 242 & -4.214 & 6.994 & 0.472 & 0.366 & -0.711 \\
 & & & & 6 & 728 & -4.925 & 7.770 & 0.776 & 0.322 & -0.691 \\
 & & & & 7 & 2186 & -5.626 & 8.384 & 0.615 & 0.351 & -0.748 \\
\addlinespace
3 & $u^2+2$ & 1.703 & 0.500 & 1 & 2 & -0.801 & 1.000 & -- & 0.332 & 0.030 \\
 & & & & 2 & 8 & -1.803 & 1.699 & 0.700 & 0.295 & -0.309 \\
 & & & & 3 & 26 & -2.019 & 2.273 & 0.574 & 0.311 & 0.065 \\
 & & & & 4 & 80 & -2.827 & 2.924 & 0.651 & 0.282 & -0.121 \\
 & & & & 5 & 242 & -3.372 & 3.476 & 0.551 & 0.294 & -0.103 \\
 & & & & 6 & 728 & -3.949 & 4.061 & 0.585 & 0.297 & -0.091 \\
 & & & & 7 & 2186 & -4.528 & 4.661 & 0.600 & 0.292 & -0.075 \\
\addlinespace
5 & $u$ & 0.956 & 0.721 & 1 & 4 & -1.768 & 1.209 & -- & 0.366 & -0.914 \\
 & & & & 2 & 24 & -2.958 & 2.803 & 1.594 & 0.237 & -0.638 \\
 & & & & 3 & 124 & -3.770 & 3.728 & 0.925 & 0.262 & -0.501 \\
 & & & & 4 & 624 & -4.733 & 4.774 & 1.046 & 0.266 & -0.414 \\
 & & & & 5 & 3124 & -5.769 & 5.845 & 1.071 & 0.258 & -0.388 \\
\addlinespace
5 & $u^2+2$ & 1.247 & 0.701 & 1 & 4 & -0.978 & 1.197 & -- & 0.287 & -0.194 \\
 & & & & 2 & 24 & -2.250 & 2.493 & 1.296 & 0.142 & -0.316 \\
 & & & & 3 & 124 & -2.786 & 2.986 & 0.493 & 0.215 & -0.286 \\
 & & & & 4 & 624 & -3.701 & 3.857 & 0.871 & 0.202 & -0.342 \\
 & & & & 5 & 3124 & -4.465 & 4.671 & 0.813 & 0.197 & -0.298 \\
\addlinespace
7 & $u$ & 0.979 & 0.434 & 1 & 6 & -1.474 & 1.245 & -- & 0.196 & -0.468 \\
 & & & & 2 & 48 & -2.392 & 2.451 & 1.207 & 0.167 & -0.208 \\
 & & & & 3 & 342 & -3.478 & 3.455 & 1.003 & 0.169 & -0.288 \\
 & & & & 4 & 2400 & -4.439 & 4.468 & 1.013 & 0.171 & -0.235 \\
\bottomrule

\end{tabular}
\end{table}

\begin{table}[h]
\centering\footnotesize\setlength{\tabcolsep}{5pt}
\caption{The finite-support formula checked: finite sum, series, their difference, and
$\mathrm{Off}_f(N)$ from Table~\ref{tab:ff}.}\label{tab:fftail}
\begin{tabular}{rlrrrrrrr}
\toprule
$q$ & $D$ & $N$ & $M(N)$ & $\#d$ & finite sum & series & difference & $\mathrm{Off}_f(N)$ from LHS\\
\midrule
3 & $u$ & 1 & 1 & 0 & 0.0000 & 1.8322 & -1.8322 & -1.8322 \\
 & & 2 & 3 & 6 & 1.5993 & 2.2980 & -0.6988 & -0.6988 \\
 & & 3 & 6 & 199 & 1.9465 & 3.5913 & -1.6448 & -1.6448 \\
 & & 4 & 9 & 4726 & 1.7928 & 2.9146 & -1.1218 & -1.1218 \\
 & & 5 & 12 & 112209 & 2.3993 & 3.7307 & -1.3314 & -1.3314 \\
\addlinespace
3 & $u^2+1$ & 1 & 2 & 1 & 0.0000 & 0.2367 & -0.2367 & -0.2367 \\
 & & 2 & 3 & 5 & 0.0000 & 0.7102 & -0.7102 & -0.7102 \\
 & & 3 & 6 & 146 & 1.2164 & 2.3082 & -1.0918 & -1.0918 \\
 & & 4 & 9 & 3500 & 0.9700 & 1.5557 & -0.5857 & -0.5857 \\
 & & 5 & 12 & 83301 & 1.1099 & 1.8208 & -0.7110 & -0.7110 \\
\addlinespace
3 & $u^2+2$ & 1 & 2 & 2 & 0.1999 & 0.1696 & 0.0303 & 0.0303 \\
 & & 2 & 3 & 6 & 0.0000 & 0.3088 & -0.3088 & -0.3088 \\
 & & 3 & 6 & 199 & 0.4179 & 0.3528 & 0.0650 & 0.0650 \\
 & & 4 & 9 & 4726 & 0.2862 & 0.4071 & -0.1209 & -0.1209 \\
 & & 5 & 12 & 112209 & 0.4283 & 0.5308 & -0.1025 & -0.1025 \\
\addlinespace
5 & $u$ & 1 & 1 & 0 & 0.0000 & 0.9141 & -0.9141 & -0.9141 \\
 & & 2 & 3 & 33 & 0.1063 & 0.7447 & -0.6384 & -0.6384 \\
 & & 3 & 6 & 4032 & 0.7041 & 1.2048 & -0.5006 & -0.5007 \\
\addlinespace
5 & $u^2+2$ & 1 & 2 & 5 & 0.6271 & 0.8215 & -0.1943 & -0.1943 \\
 & & 2 & 3 & 28 & 0.1140 & 0.4301 & -0.3160 & -0.3160 \\
 & & 3 & 6 & 3314 & 0.6868 & 0.9730 & -0.2862 & -0.2862 \\
\addlinespace
7 & $u$ & 1 & 1 & 0 & 0.0000 & 0.4679 & -0.4679 & -0.4679 \\
 & & 2 & 3 & 96 & 0.3520 & 0.5597 & -0.2077 & -0.2077 \\
 & & 3 & 6 & 29328 & 0.4140 & 0.7022 & -0.2882 & -0.2882 \\
\bottomrule

\end{tabular}
\end{table}

\section{Assessment}

Route B is the shorter path to a theorem: Theorem~\ref{thm:B1} is proved, Proposition~\ref{prop:B2} is a
standard consequence of Weil's theorem, and Proposition~\ref{prop:B3} needs one equidistribution input
that exists in the literature in the large-$q$ regime. Route A is the harder and more valuable path,
since it concerns the integers; it is a transcription of Hooley's 1963 method with multiplicative weights,
and we have identified exactly where the two estimates (Propositions~\ref{prop:A2} and~\ref{prop:A3})
enter. Neither is done here.

\bibliographystyle{plain}
\bibliography{refs}
\end{document}
